% !TeX root = ../../Book3.tex
% 纲要标题：## 附录 A　伪黎曼与 Lorentz 几何
% 纲要位置：几何学计划纲要.md，第 3351 行。
% 正文状态：未撰写；本文件目前是框架。

% 纲要内容开始
% 
% > 从不定度量出发说明哪些黎曼证明失效；不将完备性、距离和因果性混用。
% 
% * **局部结构与测地线**
%   * 不定配对及 Levi-Civita 联络
%   * 时间、空间和零向量
%   * 测地线和光锥
%   * 度量完备性的不同含义
% 
% * **因果结构与整体问题**
%   * 时间定向和因果曲线
%   * Cauchy 超曲面及全局双曲性
%   * Minkowski 与典型时空模型
%   * 奇点定理的条件和外部输入
% 
% ---
% 
% 纲要内容结束

\chapter{伪黎曼与 Lorentz 几何}
\label{g3:app:A}
\BookChapterNumber{g3:app:A}

待撰写

\section{局部结构与测地线}
\label{g3:sec:A01}
待撰写

\subsection{不定配对及 Levi-Civita 联络}
\label{g3:subsec:A01:01}
待撰写

\subsection{时间、空间和零向量}
\label{g3:subsec:A01:02}
待撰写

\subsection{测地线和光锥}
\label{g3:subsec:A01:03}
待撰写

\subsection{度量完备性的不同含义}
\label{g3:subsec:A01:04}
待撰写

\section{因果结构与整体问题}
\label{g3:sec:A02}
待撰写

\subsection{时间定向和因果曲线}
\label{g3:subsec:A02:01}
待撰写

\subsection{Cauchy 超曲面及全局双曲性}
\label{g3:subsec:A02:02}
待撰写

\subsection{Minkowski 与典型时空模型}
\label{g3:subsec:A02:03}
待撰写

\subsection{奇点定理的条件和外部输入}
\label{g3:subsec:A02:04}
待撰写
